\section{Seeding the Summarizer with Gold Quasi-Sentence States}\label{app:gold-seeding}

Section~\ref{sec:gold-seeding} summarizes the gold-seeding experiment.
This appendix specifies the protocol, documents the supervision leak
the first design contained (and the diagnostic signature that exposed
it), and reports the per-dimension grid for the corrected protocol.

\subsection{Design}

The experiment asks a budgeted local-to-global question: given a small
random sample of leaf-window states from a document, can a learned
summarizer $g$ produce a full-document compact state that a fixed,
already-trained readout $f$ maps to the document gold? The sampling
design mirrors the audit of Section~\ref{sec:estimation} (known
propensities over leaf windows), so a positive answer means the same
sampled infrastructure that certifies a tree can also seed one.

The protocol, per document: fix a trained $f$ artifact from the
$\ell = 16$ ladder (Appendix~\ref{app:qsentence-leaf-ladder}); sample
$16$ leaf windows; represent each sampled window by a state; prompt-optimize
$g$ (GEPA) to compose those states into a document-level compact state,
with a reward that mixes direct parseability against the document gold
($0.75$) and agreement of $f$ applied to $g$'s output ($0.25$). Train
on $140$ documents, evaluate on the $48$ test documents. Three
estimators are compared against the full teacher rollup:
\begin{itemize}[leftmargin=1.5em]
    \item \textbf{sample baseline}: aggregate the sampled windows'
          states directly (no $g$); this is what the sample alone is
          worth.
    \item \textbf{g\_direct}: parse the document state from $g$'s
          output.
    \item \textbf{f\_on\_g}: apply the fixed $f$ to $g$'s output; this
          is the deployment path.
\end{itemize}

\subsection{The Gold-State Leak and Its Signature}

The load-bearing design choice is \emph{what state represents a
sampled window}, controlled by \texttt{--sample-state-source}. The
first implementation used \texttt{gold\_stats}: the window's gold CMP
weighted means, injected into $g$'s prompt. The resulting runs scored
deceptively well and taught nothing: the pipeline regressed gold onto
gold, $g$'s learned behavior added noise on top of the injected
answer (it underperformed the sample baseline built from the same
gold states), and nothing in the loop ever read text. The corrected
source, \texttt{f\_states} (now the default), represents each sampled
window by the learned $f$'s own prediction for that window, so gold
enters training targets only and never the inference path.

The signature that exposed the leak generalizes to any
tree-supervision pipeline and is cheap to monitor: a widening gap
between internal agreement and external agreement. A related
exposure-bias variant appears in the scheduled-sampling control run,
where a merge trained exclusively on gold child states collapses when
fed its own generated children at evaluation (internal $r = 0.166$
while the external number stays high at $0.916$ off the gold rollup).
Both failures are invisible to any single headline metric and obvious
in the pair; this is the operational content of reporting internal and
external columns together throughout
Appendix~\ref{app:qsentence-leaf-ladder}.

\subsection{Results Under the Corrected Protocol}

Table~\ref{tab:app-gold-seeding} reports the corrected
(\texttt{f\_states}) run at $\ell = 16$, $16$ sampled windows, one
sample per document, test split ($n = 48$), Pearson against the full
teacher rollup.

\begin{table}[htbp]
    \centering
    \begin{tabular}{lccc}
    \toprule
    dimension & sample baseline & g\_direct & f\_on\_g \\
    \midrule
    all dims  & $0.966$ & $0.953$ & $0.951$ \\
    RILE      & $0.949$ & $0.955$ & $0.957$ \\
    domain\_1 & $0.749$ & $0.738$ & $0.741$ \\
    domain\_2 & $0.920$ & $0.917$ & $0.910$ \\
    domain\_3 & $0.933$ & $0.890$ & $0.898$ \\
    domain\_4 & $0.876$ & $0.859$ & $0.840$ \\
    domain\_5 & $0.937$ & $0.898$ & $0.899$ \\
    domain\_6 & $0.946$ & $0.937$ & $0.923$ \\
    domain\_7 & $0.697$ & $0.441$ & $0.447$ \\
    \bottomrule
    \end{tabular}
    \caption{Gold-seeded $g$ under the corrected \texttt{f\_states}
    protocol: Pearson vs.\ the full teacher rollup, $\ell = 16$, $16$
    sampled leaf windows per document, one sample per document, test
    split ($n=48$).}
    \label{tab:app-gold-seeding}
\end{table}

Three numbers summarize the outcome. The seeded $g$ tracks the full
rollup at all-dimension $r = 0.953$ against a sample-only baseline of
$0.966$: composing through the learned summarizer costs about one
point of correlation relative to using the raw sample, and on RILE the
composed path is slightly better than the baseline ($0.955$ vs.\
$0.949$). Agreement between the seeded $g$ and its own sample baseline
is $0.984$ all-dims, so nearly all remaining error is the sampling
budget, not the learned composition. The deployment path
\texttt{f\_on\_g} matches \texttt{g\_direct} to within $0.002$
all-dims, so reading $g$'s state through the trained readout is
lossless in practice. Domain\_7 is the one dimension where the learned
composition genuinely lags its own sample ($0.441$ vs.\ $0.697$); it
is the sparsest dimension and the same one every experiment in
Appendix~\ref{app:qsentence-leaf-ladder} flags.

\subsection{Open Arms}

Two companion runs are scoped and incomplete; both are stated here as
placeholders so their eventual numbers land in a fixed frame.

\paragraph{Scheduled-sampling A/B.} The exposure-bias fix for the
merge (train on a schedule mixing gold and self-generated children)
has a completed control arm and an untrained treatment arm. The
control (rate $0$: gold children only, whose collapse signature is
quoted above) completed in every run. The treatment (rate $1$) never
reached $g$-training: neither launch of the leaf-$8$ A/B contains a
trained $g$ artifact or post-training evaluation records, and the
leaf-$2$ run's scheduled arm never started. Resuming means running
the treatment $g$-training stage
(\texttt{scripts/run\_scheduled\_sampling\_ab\_leaf8.sh} with
\texttt{--dspy-g-scheduled-sampling-rate 1.0}) and comparing against
the saved control. \fixme{scheduled-sampling treatment arm (rate 1)
g-training and evaluation pending; table to be added on completion.}

\paragraph{Cross-runtime transfer.} Prompt programs trained with
Gemma-4 against gold quasi-sentence labels, evaluated under the
DiffusionGemma runtime, would separate ``the program learned the
task'' from ``the program learned the runtime.'' The earlier transfer
wrappers produced empty reports because they pointed at a source run
that crashed before producing $g$ artifacts; complete gold-trained
$f$+$g$ artifacts exist elsewhere, at leaves $8$ and $16$
(\texttt{outputs/manifesto\_parallel\_llm\_qsentence\_20260621\_075132/gemma4\_fixed\_leafgrid/dspy/})
and leaf $1$
(\texttt{outputs/manifesto\_qsentence\_followon\_all\_missing\_20260623\_182238/gemma4\_full\_leafgrid/dspy/leafq001/}).
The transfer script
(\texttt{scripts/run\_manifesto\_qsentence\_gemma4\_gold\_to\_dgemma.sh})
evaluates from existing artifacts via the \texttt{SOURCE\_RUN\_ROOT},
\texttt{SOURCE\_GRID}, and \texttt{LEAF\_QS} overrides, so the
evaluation is runnable without retraining at leaves $1$, $8$, and
$16$; leaves $2$ and $4$ have no gold-trained $g$ and would need
training first. \fixme{gold-to-DiffusionGemma transfer evaluation
pending at leaves 1, 8, 16; table to be added on completion.}
